% Appendix: its own sections, so the mini-frame headline shows only these.
\section[Backup]{Backup}

% ------------------------------------------------------------ loss terms
\begin{frame}[source={Liao et al.\ (2026), Nat.\ Commun.\ 17, 11232, Eqs.\ 2--15}]{A1\enspace The loss terms in full}
  \begin{columns}[T,onlytextwidth]
    \begin{column}{0.48\textwidth}
      \small
      {\scriptsize\color{udsgrey}Sobel kernels and edge term (Eqs.\ 2, 3)}
      \[ S_x=\left[\begin{array}{rrr}-1&0&1\\-2&0&2\\-1&0&1\end{array}\right],\;\;
         S_y=\left[\begin{array}{rrr}-1&-2&-1\\0&0&0\\1&2&1\end{array}\right] \]
      \[ L_{\mathrm{edge}}=\mathrm{mean}\big(|S_x\ast x-S_x\ast\hat{x}|+|S_y\ast x-S_y\ast\hat{x}|\big) \]
      {\scriptsize\color{udsgrey}structural similarity (Eqs.\ 8, 9)}
      \[ \mathrm{SSIM}(x,\hat{x})=\frac{(2\mu_x\mu_{\hat{x}}+C_1)(2\sigma_{x\hat{x}}+C_2)}
         {(\mu_x^2+\mu_{\hat{x}}^2+C_1)(\sigma_x^2+\sigma_{\hat{x}}^2+C_2)} \]
      \[ L_{\mathrm{SSIM}}=1-\mathrm{SSIM}(x,\hat{x}) \]
    \end{column}
    \begin{column}{0.48\textwidth}
      \small
      {\scriptsize\color{udsgrey}VAE, with $x_V$ the reconstruction of $x$ (Eq.\ 11)}
      \[ L_{\mathrm{VAE}}=L_{\mathrm{pix}}+L_{\mathrm{edge}}+L_{\mathrm{SSIM}}+L_{\mathrm{KL}}\big(q(z|x),p(z)\big) \]
      {\scriptsize\color{udsgrey}diffusion: forward process (Eq.\ 13)}
      \[ q(x_t|x)=\mathcal{N}\big(x_t;\sqrt{\bar{\alpha}_t}\,x,(1-\bar{\alpha}_t)\mathbf{I}\big) \]
      {\scriptsize\color{udsgrey}diffusion: loss on the predicted noise $\epsilon_\theta(x_t,t,x_V)$ (Eq.\ 15)}
      \[ L_{\mathrm{DDPM}}=L_{\mathrm{pix}}(\epsilon_\theta,\tilde{\epsilon})+L_{\mathrm{SSIM}}(\epsilon_\theta,\tilde{\epsilon})
         +L_{\mathrm{edge}}(\epsilon_\theta,\tilde{\epsilon}) \]
    \end{column}
  \end{columns}
  \note{\textbf{If asked}\par
    The pixel, edge and SSIM terms are added without weights, in the VAE and in the diffusion model.\par
    In the diffusion model they compare the predicted with the true noise, so their link to grain
      boundaries is indirect.\par
    Eq.\ 14 as printed gives the predicted noise as the mean of the reverse step; standard DDPM sampling
      derives that mean from the predicted noise.\par
    SSIM uses its three factors (luminance, contrast, structure) with equal exponents of one, which
      gives Eq.\ 8; $C_1$ and $C_2$ are small constants.}
\end{frame}

% ----------------------------------------------------------------- the data
\begin{frame}[source={Liao et al.\ (2026), Nat.\ Commun.\ 17, 11232}]{A2\enspace The 25 processing routes}
  \scriptsize\renewcommand\arraystretch{1.08}\setlength\tabcolsep{6pt}
  \def\alloyhead{\toprule
      & \multicolumn{2}{c}{annealing} & rolling & \multicolumn{2}{c}{measured}\\
    ID & °C & h & \% & MPa & \%\\
    \midrule}
  \hfill
  \begin{tabular}[t]{@{}rrrrrr@{}}
    \alloyhead
    \alloytableA
    \bottomrule
  \end{tabular}\hfill\hfill
  \begin{tabular}[t]{@{}rrrrrr@{}}
    \alloyhead
    \alloytableB
    \bottomrule
  \end{tabular}\hfill\null
  \note{\textbf{If asked}\par
    The table of the code repository: each variant's processing route, hot rolling at 900 °C and then
      annealing, with its measured yield strength and elongation.\par
    Every pair of levels occurs once: an orthogonal design.\par
    Alloy 10 is the measured alloy closest in strength to candidate 1.\par
    The paper's 20/5 split is random and not listed; from Fig.\ 4b, the test alloys lie near 375, 445,
      470, 795 and 845 MPa.}
\end{frame}

\section[References]{References}

\begin{frame}{A3\enspace References}
  \begin{columns}[T,onlytextwidth]
  \begin{column}{0.53\textwidth}
  \scriptsize\setlength\parskip{0.6mm}\raggedright
  \textbf{The paper and its code}\par
  W.~Liao, K.~Li, B.~Tang, J.~Fan, J.~Wang, X.~Xue, J.~Li, R.~Yuan (2026). \emph{Generative design of high-fidelity
  microstructures using physics-aware machine learning}. Nature Communications 17, 11232.
  \url{https://doi.org/10.1038/s41467-026-78118-3}, CC~BY-NC-ND~4.0.\par
  W.~Liao, R.~Yuan (2026). \emph{microstructures\_design}, code, trained models and \fn{dataset.csv}, commit 25f6017,
  MIT licence, read 4~Nov~2026.\\ \url{https://github.com/nwpuai4msegroup/microstructures_design};\\ archived at
  \url{https://doi.org/10.5281/zenodo.21897025}.\par
  \textbf{Materials background}\par
  E.~Ma, T.~Zhu (2017). \emph{Towards strength--ductility synergy through the design of heterogeneous nanostructures
  in metals}. Materials Today 20(6), 323--331.\par
  \end{column}
  \begin{column}{0.43\textwidth}
  \scriptsize\setlength\parskip{0.4mm}\raggedright
  \textbf{Generative models}\par
  D.~P.~Kingma, M.~Welling (2013). \emph{Auto-encoding variational Bayes}. arXiv:1312.6114; ICLR 2014.\par
  I.~Goodfellow et al.\ (2014). \emph{Generative adversarial nets}. NeurIPS~27.\par
  J.~Sohl-Dickstein, E.~Weiss, N.~Maheswaranathan, S.~Ganguli (2015). \emph{Deep unsupervised learning using
  nonequilibrium thermodynamics}. ICML 2015, PMLR 37, 2256--2265.\par
  J.~Ho, A.~Jain, P.~Abbeel (2020). \emph{Denoising diffusion probabilistic models}. NeurIPS~33.\par
  A.~Ramesh et al.\ (2021). \emph{Zero-shot text-to-image generation}. ICML 2021, PMLR 139, 8821--8831.\par
  R.~Rombach, A.~Blattmann, D.~Lorenz, P.~Esser, B.~Ommer (2022). \emph{High-resolution image synthesis with latent
  diffusion models}. CVPR 2022, 10684--10695.\par
  K.~Pandey, A.~Mukherjee, P.~Rai, A.~Kumar (2022). \emph{DiffuseVAE: efficient, controllable and high-fidelity
  generation from low-dimensional latents}. Transactions on Machine Learning Research.\par
  \end{column}
  \end{columns}
\end{frame}

\begin{frame}{A4\enspace References}
  \begin{columns}[T,onlytextwidth]
  \begin{column}{0.48\textwidth}
  \scriptsize\setlength\parskip{0.6mm}\raggedright
  \textbf{Image processing}\par
  I.~Sobel, G.~Feldman (1968). \emph{A 3\,×\,3 isotropic gradient operator for image processing}. Talk at the Stanford
  Artificial Intelligence Project.\par
  Z.~Wang, A.~C.~Bovik, H.~R.~Sheikh, E.~P.~Simoncelli (2004). \emph{Image quality assessment: from error visibility
  to structural similarity}. IEEE Transactions on Image Processing 13(4), 600--612.\par
  O.~Ronneberger, P.~Fischer, T.~Brox (2015). \emph{U-Net: convolutional networks for biomedical image
  segmentation}. MICCAI 2015, 234--241.\par
  \textbf{Optimization}\par
  V.~Pareto (1906). \emph{Manuale di economia politica}. Società Editrice Libraria, Milan.\par
  J.~H.~Holland (1975). \emph{Adaptation in Natural and Artificial Systems}. University of Michigan Press.\par
  K.~Deb, A.~Pratap, S.~Agarwal, T.~Meyarivan (2002). \emph{A fast and elitist multiobjective genetic algorithm:
  NSGA-II}. IEEE Transactions on Evolutionary Computation 6(2), 182--197.\par
  \end{column}
  \begin{column}{0.48\textwidth}
  \scriptsize\setlength\parskip{0.6mm}\raggedright
  \textbf{Statistics}\par
  K.~Pearson (1901). \emph{On lines and planes of closest fit to systems of points in space}. Philosophical
  Magazine 2(11), 559--572.\par
  L.~Breiman, J.~H.~Friedman, R.~A.~Olshen, C.~J.~Stone (1984). \emph{Classification and Regression Trees}.
  Wadsworth.\par
  J.~H.~Friedman (2001). \emph{Greedy function approximation: a gradient boosting machine}. Annals of Statistics
  29(5), 1189--1232.\par
  \end{column}
  \end{columns}
\end{frame}
